% Standalone problem document, generated from the adjacent README.md.
% Compile with XeLaTeX. Mathematical content is maintained in Markdown.
\documentclass[11pt,a4paper]{article}
\usepackage[fontsize=10.5pt]{scrextend}
\usepackage[margin=22mm,headheight=15pt,headsep=9mm,footskip=11mm]{geometry}
\usepackage{fontspec}
\setmainfont{texgyrepagella}[Extension=.otf,UprightFont=*-regular,BoldFont=*-bold,ItalicFont=*-italic,BoldItalicFont=*-bolditalic]
\setsansfont{texgyreheros}[Extension=.otf,UprightFont=*-regular,BoldFont=*-bold,ItalicFont=*-italic,BoldItalicFont=*-bolditalic]
\setmonofont{texgyrecursor}[Extension=.otf,UprightFont=*-regular,BoldFont=*-bold,ItalicFont=*-italic,BoldItalicFont=*-bolditalic,Scale=MatchLowercase]
\usepackage{amsmath,amssymb}
\usepackage{unicode-math}
\setmathfont{texgyrepagella-math.otf}
\usepackage{xcolor,microtype,enumitem,needspace,fancyhdr}
\usepackage{bookmark}
\definecolor{ink}{HTML}{17344B}
\definecolor{accent}{HTML}{197D80}
\definecolor{muted}{HTML}{596773}
\hypersetup{colorlinks=true,linkcolor=accent,urlcolor=accent,pdftitle={MD-04 - The
Beck--Fiala discrepancy
conjecture},pdfauthor={Open Problems in Numerical Linear Algebra}}
\urlstyle{same}
\setlength{\parindent}{0pt}
\setlength{\parskip}{5pt plus 1pt minus 1pt}
\linespread{1.04}
\setlength{\emergencystretch}{3em}
\widowpenalty=10000
\clubpenalty=10000
\displaywidowpenalty=10000
\allowdisplaybreaks[1]
\setlist{nosep,leftmargin=1.5em}
\providecommand{\tightlist}{\setlength{\itemsep}{0pt}\setlength{\parskip}{0pt}}
\providecommand{\pandocbounded}[1]{#1}
\pagestyle{fancy}
\fancyhf{}
\fancyhead[L]{\sffamily\footnotesize\color{muted}OPEN PROBLEMS IN NUMERICAL LINEAR ALGEBRA}
\fancyhead[R]{\sffamily\bfseries\footnotesize\color{accent}MD-04}
\fancyfoot[L]{\parbox[b]{0.9\textwidth}{\sffamily\fontsize{8}{10}\selectfont\color{muted}Editorial ratings; a bounded search cannot certify that a problem remains unsolved.\\Literature check: 2026-09-10}}
\fancyfoot[R]{\sffamily\footnotesize\color{muted}\thepage}
\renewcommand{\headrulewidth}{0.3pt}
\renewcommand{\headrule}{\hbox to\headwidth{\color{accent}\leaders\hrule height \headrulewidth\hfill}}
\makeatletter
\renewcommand{\section}{\Needspace{5\baselineskip}\@startsection{section}{1}{0pt}{12pt plus 2pt minus 2pt}{4pt}{\sffamily\bfseries\large\color{ink}}}
\renewcommand{\subsection}{\Needspace{4\baselineskip}\@startsection{subsection}{2}{0pt}{9pt plus 1pt minus 1pt}{3pt}{\sffamily\bfseries\normalsize\color{ink}}}
\makeatother
\setcounter{secnumdepth}{0}
\begin{document}
{\sffamily\small\color{muted}Matrix discrepancy and optimization\par}
\vspace{3pt}
{\raggedright\hyphenpenalty=10000\exhyphenpenalty=10000\sffamily\bfseries\fontsize{21}{24}\selectfont\color{ink}The
Beck--Fiala discrepancy conjecture\par}
\vspace{7pt}
{\sffamily\small\textbf{Difficulty:} extreme\quad\textbf{Importance:} broadly
interesting\par}
{\sffamily\small\bfseries\color{accent}Status: Partially resolved\par}
\vspace{4pt}
{\color{accent}\hrule height 0.8pt}
\vspace{6pt}
\textbf{Provenance:} source-stated conjecture.

\textbf{Rating rationale:} The unrestricted square-root sparsity law is
a central discrepancy conjecture with broad consequences for
combinatorics, rounding and optimization.

Does a finite universal constant \(C>0\) exist such that the following
holds for every pair of positive integers \(m,n\), every integer \(t\)
with \(1\leq t\leq m\), and every matrix
\(A=(a_{ij})\in\{0,1\}^{m\times n}\)? If \[
\sum_{i=1}^m a_{ij}\leq t\qquad(1\leq j\leq n),
\] then some signing \(x\in\{-1,1\}^n\) satisfies \[
\|Ax\|_\infty\leq C\sqrt t.
\] The constant must be independent of \(m,n,t\). All columns are
available when choosing the signing. The zero-sparsity case is trivial
and excluded from the quantifiers.

The matrix is the incidence matrix of a set system: each element belongs
to at most \(t\) sets, and the signing should balance every set. The
problem asks whether column sparsity alone controls the simultaneous row
error at its conjectured square-root scale. It is a special case of the
Komlós conjecture after column normalization, but an independent
classical source-stated conjecture for incidence matrices.

\subsection{References}\label{references}

\begin{enumerate}
\def\labelenumi{\arabic{enumi}.}
\tightlist
\item
  D. J. Altschuler and K. Tikhomirov, \emph{Online Beck--Fiala Down to
  Logarithmic Sparsity}, arXiv:2607.14238 (2026), abstract and
  introduction: the unrestricted conjecture and the proved sparsity
  regime. \href{https://arxiv.org/html/2607.14238v1}{Paper}.
\item
  N. Bansal and H. Jiang, \emph{Decoupling via Affine
  Spectral-Independence: Beck-Fiala and Komlós Bounds Beyond
  Banaszczyk}, arXiv:2508.03961v2 (2025), abstract and introductory
  results. \href{https://arxiv.org/abs/2508.03961}{Paper}.
\item
  J. Beck and T. Fiala, \emph{``Integer-making'' theorems}, Discrete
  Applied Mathematics 3(1) (1981), pp.~1--8, the original
  bounded-column-sparsity discrepancy results.
  \href{https://doi.org/10.1016/0166-218X(81)90022-6}{Article}.
\end{enumerate}

\subsection{Status check --- 2026-09-10}\label{status-check-2026-09-10}

checked 2607.14238v1 (2026-07-15), 2508.03961v2 (2025-09-09), and
searches ``Beck Fiala conjecture solved 2026 logarithmic sparsity'' and
``Online Beck Fiala Down to Logarithmic Sparsity''. The 2026 work
reaches sparsities \(t\geq(\log n)^{1+o(1)}\) in its notation translated
to \(n\) columns; it does not establish the assertion for all
sparsities. Online lower bounds concern a more restrictive information
model. No full resolution or withdrawal was found.

\textbf{Audit update (2026-09-10):} Rechecked the July 2026
online/offline paper and searched for unrestricted Beck--Fiala results.
It proves the offline target in a substantial sparsity regime, which is
part of the displayed family, but leaves smaller sparsities unresolved.
This is a bounded literature check, not a proof that no solution exists.
\end{document}
